\documentclass[12pt,a4paper]{article}

% Encoding and fonts
\usepackage[utf8]{inputenc}
\usepackage[T1]{fontenc}
\usepackage{mathpazo} % Palatino font — BKI requires Palatino for diacritics

% Layout
\usepackage[margin=2.5cm]{geometry}
\usepackage{setspace}
\doublespacing
\usepackage{parskip}

% Tables
\usepackage{booktabs}
\usepackage{longtable}
\usepackage{array}
\usepackage{tabularx}

% Graphics
\usepackage{graphicx}
\usepackage{float}

% References and links
\usepackage[hidelinks]{hyperref}
\usepackage{natbib}
\setcitestyle{open={(},close={)},citesep={;},aysep={},notesep={:}}
% No \bibliographystyle — manual bibliography in BKI format

% Misc
\usepackage{enumitem}
\usepackage{textcomp}
\usepackage{amssymb}

% BKI heading format: First Level = Bold, Second Level = Bold Italic
\usepackage{titlesec}
\titleformat{\section}{\normalfont\large\bfseries}{}{0em}{}
\titleformat{\subsection}{\normalfont\normalsize\bfseries\itshape}{}{0em}{}
\titleformat{\subsubsection}[runin]{\normalfont\normalsize}{}{1em}{}

% Title
\title{\textbf{The Volcanic Ritual Clock: Taphonomic Calibration of Javanese Mortuary Intervals and Their Pre-Indic Austronesian Origin}}

% ANONYMIZED FOR DOUBLE-BLIND REVIEW
% Original: M.\,N.\ Amien, Lab Data Sains, Universitas Bhinneka Nusantara, Malang
\author{[Anonymised for review]}

\date{}

\begin{document}

\maketitle

\begin{abstract}
The Javanese \textit{slametan} death ritual follows communal commemorations at 3, 7, 40, 100, and 1000 days after death. This numerical sequence has no source in Hindu, Buddhist, Islamic, or Christian scripture; four successive religious waves modified the verbal content while preserving the temporal structure. The intervals correspond to forensic decomposition stages in acidic volcanic soil (pH 4.5--5.5), with the terminal 1000-day mark matching bone dissolution in tropical Andosol ($p = 0.008$, permutation test). The underlying belief structure---death as gradual process, ritual efficacy paralleling bodily decay---is pan-Austronesian, shared by 30 of 137 cultures in the Pulotu database. We propose that Javanese mortuary timing was empirically calibrated to observable decomposition in volcanic landscapes, and that the same volcanism that calibrates this ritual clock also buries the archaeological evidence of the people who kept it.

\medskip
\noindent\textbf{Keywords:} mortuary ritual, taphonomy, slametan, Austronesian, volcanic soil, Java
\end{abstract}

[Figure 0 about here: The slametan death cycle --- what generated these numbers?]

\bigskip

%% ===================================================================
%% SECTION 1: INTRODUCTION
%% ===================================================================
\section{Introduction}

\subsection{The Problem}

Death rituals are among the most conservative elements of any cultural repertoire. The anthropological literature has long recognised that mortuary practices persist across religious conversions, political upheavals, and economic transformations with a tenacity that other social institutions rarely match \citep{bloch_parry_1982, metcalf_huntington_1991}. This conservatism is conventionally explained in terms of the emotional weight of bereavement, the authority of tradition in moments of crisis, and the social functions that mortuary ritual performs for the living. What has received far less attention is the possibility that mortuary \textit{timing}---the specific intervals at which commemorative rites are performed---may encode empirical observation of a physical process: the decomposition of the human body in local soil.

This paper addresses that possibility in a particular landscape: the volcanic plains of Java.

In volcanic environments, two taphonomic processes converge on the material correlates of death. The first is chemical. Volcanic soils---classified pedologically as Andosols---are characteristically acidic, with pH values in Java typically ranging from 4.5 to 5.5 (ISRIC SoilGrids v2.0). At this acidity, the mineral component of bone, hydroxyapatite (Ca\textsubscript{10}(PO\textsubscript{4})\textsubscript{6}(OH)\textsubscript{2}), enters a regime of active dissolution. Experimental studies have demonstrated measurable degradation of buried bone within weeks at pH values below 5.0 \citep{oghenemavwe_2022}, and forensic case reports from acidic environments document near-complete mineral loss in remains buried for only a few years \citep{kemp_2016}. The well-known `sand silhouettes' at Sutton Hoo---where acidic sandy soil reduced inhumations to stained outlines---illustrate the endpoint of this process. In the tropics, where temperatures of 26--28\textdegree C and high moisture accelerate both microbial activity and chemical weathering, the trajectory from intact burial to physical disappearance is compressed still further.

The second process is depositional. Java's volcanic arc produces tephra at rates that, in the alluvial plains surrounding major stratovolcanoes, translate to measurable sedimentation. \citet{newhall_2000} documented that the ninth-century Candi Sambisari near Yogyakarta was buried under approximately 6.5~m of deposits---40~cm of primary volcanic ash overlain by 6~m of subsequent fluvial sedimentation---discovered by a farmer in 1966 \citep{newhall_2000}. Nearby Candi Kedulan was found under approximately 7~m of 14 distinct fluvio-volcanic layers \citep{newhall_2000}. The Dwarapala guardian statues of Singhasari (East Java), carved circa 1268~CE and standing 370~cm tall, were found by Engelhard in 1803 with 185~cm buried beneath volcanic deposits---implying approximately 3.5~mm per year over 535 years. At composite rates of 3.5--7~mm per year (mixing catastrophic events with inter-eruptive reworking), an archaeological site is buried beneath several metres of deposit within a millennium---well beyond the depth of casual discovery or agricultural disturbance. Open-air sites in volcanic Java are not merely poorly preserved; they are systematically rendered invisible to surface survey.

These two processes operate at different timescales---years for decomposition, centuries for burial---but they share a common driver: volcanism. The same tephra that weathers into acidic soil also accumulates as the sedimentary matrix that entombs sites. This convergence raises a question that, to the best of our knowledge, has not been posed in the existing literature: could the timing of mortuary ritual in volcanic landscapes reflect generations of empirical observation of how bodies decompose in that specific soil?

\subsection{The Slametan Cycle}

The Javanese \textit{slametan} (also \textit{selamatan}, \textit{kenduren}) is a communal ritual meal that marks transitions of many kinds---birth, circumcision, house-building, harvest, calendrical events, and, most elaborately, death \citep{geertz_1960, beatty_1999}. In its mortuary variant, the slametan follows a structured sequence of commemorations after the event of death. The intervals are specific and consistent across regions, social classes, and---remarkably---religious affiliations. Muslim, Catholic, Protestant (GKJ), and Kejawen (Javanese traditional religion) communities all observe some version of the same sequence, differing in the verbal content of prayers but not in the timing or social structure of the gatherings.

The standard mortuary sequence runs as follows: \textit{nelung dina} (third day), \textit{mitung dina} (seventh day), \textit{matang puluh} (fortieth day), \textit{nyatus} (hundredth day), \textit{mendhak sepisan} (first anniversary), \textit{mendhak pindo} (second anniversary), and \textit{nyewu} (thousandth day). Each successive commemoration is typically larger than the last, involving more guests, longer recitations, and greater expenditure. The sequence builds toward \textit{nyewu} as its culmination---the final ceremony, after which no further commemorations are held and the deceased is considered to have completed the transition from the world of the living \citep{geertz_1960, koentjaraningrat_1994, hendrajaya_2020}.

What makes this sequence anthropologically distinctive is the explicit connection drawn by Javanese tradition between the timing of each ceremony and the physical state of the body. This is not an inference by the present authors; it is stated in Javanese sources themselves. At \textit{matang puluh} (forty days), the ceremony is understood to `perfect' (\textit{menyempurnakan}) specific bodily components: `blood, flesh, marrow, innards, nails, hair, bones, and muscle' (\textit{darah, daging, sungsum, jeroan, kuku, rambut, tulang, dan otot})---a systematic inventory of tissues, not a vague metaphor \citep[437]{hendrajaya_2020}. At \textit{mendhak pindo} (approximately two years), the tradition holds that `the body has almost dissolved, only bones remain' (\textit{jenazah sudah hampir luluh, tinggal tulang saja})---a claim independently documented in two separate academic sources \citep[439]{hendrajaya_2020, sholihah_2023}. At \textit{nyewu} (one thousand days), Geertz recorded that the body is believed to have `completely decayed to dust' \citep[ch.~6]{geertz_1960}.

The parallel between spiritual departure and physical decomposition runs through the entire sequence. At three days, the spirit is believed to remain in the house. At seven days, it begins to depart. At forty days, it leaves the immediate surroundings. At one thousand days, it will `no longer return to the family' \citep[440]{hendrajaya_2020}. The soul's progressive journey outward from the domestic sphere mirrors, step by step, the progressive dissolution of the body beneath the earth. The critical observation---and the one on which this paper's argument turns---is that the ritual's endpoint is defined by a physical criterion, not a theological one. \textit{Nyewu} is not the day on which a doctrine declares the soul complete; it is the day on which, in volcanic soil, the body can no longer be found.

\subsection{Prior Work}

The theoretical framework connecting physical decomposition to mortuary ritual has a distinguished pedigree. Robert Hertz, in his 1907 essay \textit{Contribution \`{a} une \'{e}tude sur la repr\'{e}sentation collective de la mort}, proposed that in societies practising secondary burial, the period between death and final ceremony corresponds to the time required for the body to reach the `dry' state---stripped of flesh, reduced to clean bone \citep{hertz_1907}. For Hertz, the liminal status of the corpse---neither fully alive nor fully dead---was mirrored in the liminal status of the soul, which could not join the ancestors until the body had completed its transformation.

\citet{bloch_parry_1982} extended this analysis cross-culturally, demonstrating that the association between bodily decomposition and spiritual transition recurs across widely separated societies. \citet{metcalf_huntington_1991} documented the specific Austronesian manifestation, showing that secondary burial and multi-stage mortuary ritual are pervasive across the Austronesian-speaking world, from Borneo to Madagascar, from the Philippines to Polynesia.

Within Javanese studies, \citet{beatty_1999} offered a multivocal analysis of the slametan; \citet{woodward_1989} situated it within Javanese Islam. More recently, \citet{hendrajaya_2020} provided the most detailed published account of the body-state beliefs associated with each interval. \citet[149]{aizid_2015} stated explicitly that the slametan kematian `existed before Hinduism-Buddhism came to Indonesia', with successive religious waves modifying only the verbal content while leaving the timing intact.

What is absent from this literature is any attempt to test the empirical accuracy of the tradition's claims about decomposition. Hertz theorised that decomposition and ritual are linked; the Javanese sources assert specific timings for specific bodily transformations; forensic taphonomy has produced quantitative data on decomposition rates in various soil environments. Yet no previous study has brought these three bodies of knowledge into contact.

\subsection{This Paper's Contribution}

We propose what we term the `Volcanic Ritual Clock' hypothesis: that the Javanese slametan mortuary sequence preserves a set of empirically calibrated intervals corresponding to the observable stages of human decomposition in Andosol, and that this calibration reflects generations of accumulated observation in volcanic landscapes. We do not claim that the slametan is `nothing but' a decomposition calendar---its social, emotional, and cosmological dimensions are real and irreducible. We claim, rather, that the physical process of decomposition provided the empirical scaffold upon which those richer meanings were built, and that this scaffold is still legible in the tradition's own language about the body.

The argument draws on three independent lines of evidence. The first is epigraphic: analysis of the DHARMA corpus of 268 Nusantaran inscriptions reveals a persistent pre-Indic cosmological vocabulary. The second is comparative: the Pulotu database of 137 Austronesian cultures identifies at least 30 sharing the `full mortuary package' of deified ancestors, active ancestral spirits, and post-death ritual efficacy. The third is taphonomic: forensic decomposition literature provides a predicted timeline for body dissolution in tropical acidic soil that corresponds systematically to the slametan intervals.

These three lines of evidence converge on a synthesis we frame as the Hypothetical Taphonomic Overlay Model, or H-TOM. The model identifies two timescales linked by a single volcanic driver. At the ritual timescale---approximately one thousand days---acidic volcanic soil reduces a buried body to physical invisibility. At the geological timescale---approximately one thousand years---volcanic sedimentation buries the archaeological sites of the communities that performed these rituals. The environment that calibrates the ritual clock also destroys the material evidence of the people who kept it.


%% ===================================================================
%% SECTION 2: SOURCE EXCLUSION
%% ===================================================================
\section{Source Exclusion: The Slametan Sequence Is Uniquely Javanese}

The Javanese death slametan follows a fixed sequence of communal ritual meals held at 3, 7, 40, 100, and 1000 days after death, with additional commemorations at the first and second anniversaries \citep{geertz_1960, koentjaraningrat_1994}. The sequence is known locally as \textit{nelung dina} (3), \textit{mitung dina} (7), \textit{matang puluh} (40), \textit{nyatus} (100), \textit{mendhak sepisan} (1 year), \textit{mendhak pindo} (2 years), and \textit{nyewu} (1000) \citep[433--435]{hendrajaya_2020}. Because Javanese culture has been shaped by four successive waves of imported religion---Hindu-Buddhist influence from approximately the 4th century CE (earliest inscriptions at Kutai and Tarumanagara), Islam from the 13th century (Samudra Pasai) with mass conversion in the 15th--16th centuries, and Christianity from the 16th century (Portuguese missions, later Dutch Reformed and Catholic)---the immediate question is whether these intervals derive from any of these imported religious traditions.

\subsection{Hindu Sources}

The Hindu mortuary rite (\textit{antyeshti}) and its associated post-death observances (\textit{shraddha}) follow a 10--13 day cycle. The \textit{sapindikarana} ceremony, which formally incorporates the deceased into the company of ancestors (\textit{pitrs}), is performed on the 12th or 13th day after death \citep{knipe_1977, parry_1994}. Subsequent annual observances occur during \textit{Pitru Paksha}.

Crucially, the Hindu system contains no 40-day, 100-day, or 1000-day interval. The shraddha cycle operates on a base-10/12/13 arithmetic that bears no structural resemblance to the slametan sequence.

\begin{table}[H]
\centering
\caption{Hindu vs.\ Javanese Mortuary Intervals}
\label{tab:hindu_javanese}
\begin{tabular}{lll}
\toprule
\textbf{Hindu (shraddha)} & \textbf{Javanese (slametan)} & \textbf{Overlap?} \\
\midrule
Day 1--3: funeral rites & Day 0: geblag (burial) & Trivial \\
Day 10--13: sapindikarana & Day 3: nelung dina & NO \\
--- & Day 7: mitung dina & NO \\
--- & Day 40: matang puluh & NO \\
--- & Day 100: nyatus & NO \\
Annual: Pitru Paksha & Day 365: mendhak sepisan & Different system \\
--- & Day 1000: nyewu & NO \\
\bottomrule
\end{tabular}
\end{table}

\subsection{Buddhist Sources}

The Buddhist concept of the intermediate state (\textit{antarabhava} or \textit{bardo}) provides a 49-day framework (7~$\times$~7 = 49 days) for the soul's transition between death and rebirth \citep{cuevas_2003}. Some East Asian Buddhist traditions extend this with a 100-day observance, though scholars have argued that this reflects Confucian rather than Buddhist influence \citep{teiser_1988}.

The Buddhist system thus overlaps with the slametan at two points: the 7-day unit and the 100-day mark. However, it diverges fundamentally: the Buddhist 7-day cycle repeats seven times (days 7, 14, 21, 28, 35, 42, 49), whereas the Javanese system uses 7 as a single marker followed by 40. There is no 40-day interval in Buddhist doctrine. Most importantly, there is no 1000-day observance in any Buddhist scriptural tradition.

\subsection{Islamic Sources}

The Islamic mourning period (\textit{hidad}) is limited to three days, with an exception for widows (four months and ten days, per Quran 2:234). The practices of communal prayer gatherings at 3, 7, 40, 100, and 1000 days are explicitly classified as \textit{bid'ah} (innovation) by Islamic jurisprudence \citep[155--157]{woodward_1989}.

Indonesian Islamic scholars have consistently distinguished these practices as \textit{adat} (local custom) rather than \textit{syariat} (religious law). \citet[441]{hendrajaya_2020} document this process explicitly: `In the old belief, the ceremony was performed by making offerings to supernatural forces. Islam flexibly gave new colour to those ceremonies.' The key insight is one of substitution, not invention: Islam replaced the verbal content of pre-existing ceremonies while preserving the temporal structure intact.

This interpretation is supported by \citet[149]{aizid_2015}, who states directly: `The origin of the death slametan existed before Hinduism-Buddhism came to Indonesia.'

\subsection{Christian Sources}

Christianity arrived in Java significantly later than the other three traditions, with Portuguese Catholic missions in the 16th century and Dutch Reformed influence from the 17th century onward. Javanese Christian communities---both Catholic (concentrated around Muntilan and Yogyakarta) and Protestant (Gereja Kristen Jawa, GKJ)---widely practise the slametan death cycle with the same intervals. Catholic parishes in Central Java incorporate the slametan into their pastoral practice through what theologians of inculturation term `contextualisation': the prayers are Christian, but the temporal structure is inherited unchanged from the pre-Christian community \citep{aritonang_steenbrink_2008}.

The Christian case is particularly instructive for the source exclusion argument. Christianity has no indigenous tradition of commemorations at 3, 7, 40, 100, or 1000 days after death. The fact that Javanese Christians adopted the slametan intervals wholesale upon conversion---just as Javanese Muslims had done centuries earlier---confirms that the temporal sequence belongs to the Javanese substrate, not to any imported superstratum. Four successive religious waves have each modified the verbal content of the ceremony while preserving its timing intact.

\subsection{Verdict: A Uniquely Javanese Sequence}

The complete 3--7--40--100--1000 day sequence cannot be derived from any of the four imported religious traditions:

\begin{itemize}[nosep]
\item \textbf{Hinduism} (from $\sim$4th c.\ CE) provides none of the intervals beyond the trivial day-3 overlap.
\item \textbf{Buddhism} (from $\sim$4th c.\ CE) provides 7 and possibly 100, but not 40 or 1000.
\item \textbf{Islam} (from $\sim$13th c.\ CE) explicitly classifies the intervals as non-Islamic (\textit{bid'ah}) and recognises them as pre-existing local custom.
\item \textbf{Christianity} (from $\sim$16th c.\ CE) has no tradition of these intervals; Javanese Christians adopted them from the existing substrate.
\end{itemize}

The most parsimonious explanation is that the numerical sequence as a system is indigenous to Java, predating all three imported traditions. We note that individual numbers within the sequence---particularly 40---do appear in other mortuary traditions: Eastern Orthodox Christianity observes commemorations at 3, 9, and 40 days; some African traditions hold a second funeral at approximately 40 days; and ancient Egyptian mummification required 40 days of natron drying, an interval empirically calibrated to the time required to desiccate a corpse \citep{geertz_1960}. The Sufi practice of \textit{chilla} (40-day spiritual retreat) provides a plausible channel through which the number 40 may have been reinforced in Javanese usage, though the chilla is a living spiritual exercise, not a mortuary rite, and Hindu shraddha---the pre-Islamic death ritual system in Java---contains no 40-day interval.

However, the uniqueness claim advanced here concerns not individual numbers but the \textit{complete sequence}. The combination 3--7--40--100--1000 is unattested in any non-Javanese tradition. The 1000-day terminal point (\textit{nyewu}) is particularly distinctive: no other mortuary system in the comparative literature employs it. If the cross-cultural recurrence of a $\sim$40-day mortuary interval reflects independent convergence on the approximate timing of soft tissue decomposition---as the Egyptian mummification case suggests---then this convergence supports rather than undermines the taphonomic hypothesis. The question is not whether 40 appears elsewhere, but what generated the specific five-interval system and its terminal 1000-day endpoint. In the following sections, we argue that the answer lies in empirical observation of a physical process specific to volcanic soil.


%% ===================================================================
%% SECTION 3: EPIGRAPHIC EVIDENCE
%% ===================================================================
\section{Epigraphic Evidence: Pre-Indic Substrate in Old Javanese Inscriptions}

If the slametan intervals are neither Hindu, Buddhist, nor Islamic in origin, the question becomes whether an indigenous Javanese cosmological system existed prior to and independently of these imported traditions. We turn to the epigraphic record for evidence.

\subsection{Corpus and Method}

The dataset comprises 268 inscriptions from the DHARMA ERC project \citep{griffiths_2020}, a digital corpus of Southeast Asian epigraphy encoded in EpiDoc TEI-XML and released under a CC-BY 4.0 licence. The inscriptions span approximately the 5th to 15th centuries CE: Old Javanese (\textit{kawi}, $n = 155$, 58\%), Sanskrit ($n = 65$, 24\%), Old Sundanese ($n = 14$, 5\%), Old Malay ($n = 13$, 5\%), and unclassified or fragmentary texts ($n = 21$, 8\%).

We constructed a keyword ontology of 25 ritual and cosmological terms, classifying each as \textit{indic} (clear Sanskrit etymology), \textit{pre-indic} (no plausible Sanskrit source; reconstructable to Proto-Malayo-Polynesian), or \textit{ambiguous}. A term was coded as pre-Indic only if it lacked a Sanskrit etymon and possessed cognates in at least one non-Indianised Austronesian language \citep{blust_trussel_2010}.

\subsection{Results}

Pre-Indic elements were detected in 126 of 268 inscriptions (47\%). Four pre-Indic terms merit individual discussion.

\textit{Hyam/hyang} (PMP *\textit{qiang}, `spirit, divine') appeared in 116 of 268 inscriptions (43\%), making it the second most frequent ritual keyword in the entire corpus. Its distribution by language is revealing: it appears in 102/155 Old Javanese inscriptions (66\%), 6/14 Old Sundanese (43\%), 6/13 Old Malay (46\%), but in only 4/65 Sanskrit-only inscriptions (6\%). The indigenous and imported terms co-exist rather than compete: \textit{hyang} does not replace \textit{deva}, nor does \textit{deva} displace \textit{hyang}.

\textit{Manghuri} (`ancestor return') appeared in 28 of 268 inscriptions (10\%). It has no plausible Sanskrit etymon. The concept of ritually summoning ancestors has broad Austronesian parallels: the Malagasy \textit{famadihana}, the Torajan \textit{ma'nene}, and the pan-Philippine \textit{pagdiwata}.

\textit{Wuku} (210-day indigenous calendar) appeared in five inscriptions. The \textit{wuku} system has no Indian equivalent. Notably, it appears in the Old Malay inscription Dharmasraya~A from Sumatra (1286~CE), confirming the indigenous calendar circulated beyond Java.

\textit{Kabuyutan} (ancestral sacred site) appears in the Old Sundanese corpus, designating sites consecrated to ancestors.

\begin{table}[H]
\centering
\caption{Pre-Indic Elements in the DHARMA Corpus ($n = 268$)}
\label{tab:dharma}
\begin{tabular}{@{}llcccccc@{}}
\toprule
\textbf{Term} & \textbf{Etymology} & \textbf{Freq.} & \textbf{\%} & \textbf{OJ} & \textbf{Skt} & \textbf{OSn} & \textbf{OMy} \\
\midrule
hyam/hyang & PMP *qiang & 116 & 43\% & 66\% & 6\% & 43\% & 46\% \\
manghuri & Old Javanese & 28 & 10\% & --- & --- & --- & --- \\
wuku & Indigenous cal. & 5 & 2\% & --- & --- & --- & --- \\
kabuyutan & Old Sundanese & 3 & 1\% & --- & --- & \checkmark & --- \\
\bottomrule
\end{tabular}
\end{table}

\subsection{Interpretation}

The epigraphic data demonstrate that a layer of indigenous cosmology persisted throughout the period of Indianisation. Nearly half of all inscriptions (47\%) contain at least one pre-Indic element; \textit{hyang} appears in two-thirds of all Old Javanese inscriptions.

The slametan death cycle does not appear in the prasasti. This absence is expected: the inscriptions concern royal administration, not domestic mortuary practices. But the epigraphic record establishes that an indigenous cosmological vocabulary existed in formal, literate contexts; that Indic and indigenous systems co-existed through cultural layering; and that the indigenous substrate was cross-regional (Java, Sunda, Sumatra). The slametan belongs to this same substrate.


%% ===================================================================
%% SECTION 4: CROSS-AUSTRONESIAN COMPARISON
%% ===================================================================
\section{Cross-Austronesian Comparison: The Mortuary Package}

If the belief structure underlying the slametan is genuinely pre-Indic, then cognate structures should be detectable across the Austronesian world, including in populations never subject to Hindu-Buddhist or Islamic influence.

\subsection{The Pulotu Database}

Pulotu is a standardised cross-cultural database encompassing 137 Austronesian-speaking cultures \citep{watts_2015}. Three variables are of particular relevance: Q4 (ancestor deification, 0--4), Q5 (ancestral spirit activity, 0--4), and Q10 (post-death ritual efficacy, 0--4).

\subsection{The Full Mortuary Package}

We define the `full mortuary package' as the conjunction of Q4~$\geq$~2, Q5~$\geq$~2, and Q10~$\geq$~1. Thirty of 137 cultures meet all three criteria, spanning five major regions:

\begin{itemize}[nosep]
\item \textbf{Indonesia:} Southern Toraja, Karo Batak, Toba Batak, Roti, Ata Tana `Ai, Laboya
\item \textbf{Madagascar:} Merina, Tanala
\item \textbf{Philippines:} Ifugao, Visayans, Kalinga, Subanun, Tinguian
\item \textbf{Melanesia:} Trobriand Islands, Roviana, Tanna, Simbo, Berawan, Manus
\item \textbf{Polynesia:} Marquesas
\end{itemize}

The Southern Toraja and Merina of Madagascar both meet all three criteria despite a geographic separation of approximately 8,000~km and no documented post-settlement contact. In both cases, death is not an instantaneous transition but a managed process requiring sustained ritual labour.

\subsection{Structure versus Numbers}

The comparison reveals a critical distinction. The \textit{structure}---death as gradual process, post-mortem ritual efficacy---is shared across the full-package group and appears to represent an inherited Austronesian trait. The specific \textit{numbers} are locally variable:

\begin{itemize}[nosep]
\item Java: 3--7--40--100--1000 days
\item Madagascar: 3--7 year cycles (famadihana)
\item Toraja: weeks to years (wealth-dependent)
\item Philippines: secondary burial timing varies by region
\end{itemize}

This variation is precisely what the Volcanic Ritual Clock hypothesis predicts: conserved \textit{structure} combined with locally variable \textit{timing}, consistent with environmental calibration rather than cultural inheritance of fixed intervals.

\subsection{Environmental Context}

\begin{table}[H]
\centering
\caption{Mortuary Timing and Environmental Context in Three Austronesian Cultures}
\label{tab:comparison}
\begin{tabularx}{\textwidth}{@{}lXXX@{}}
\toprule
\textbf{Parameter} & \textbf{Java (Slametan)} & \textbf{Madagascar (Famadihana)} & \textbf{Toraja (Rambu Solo)} \\
\midrule
Mortuary timing & 3--7--40--100--1000 days & 3--7 year cycles & Variable (weeks--years) \\
Terminal belief & Body fully one with earth & Bones clean enough to rewrap & Soul's journey complete \\
Soil type & Volcanic andosol & Highland laterite & Volcanic highland \\
Soil pH & 4.5--5.5 & 4.7 & 5.0 \\
Calibration & Taphonomic (intervals track decomposition) & Direct observation (exhumation) & Economic, but decomposition-aware \\
\bottomrule
\end{tabularx}
\end{table}

At the aggregate level, cultures with the full mortuary package have a mean soil pH of 5.36 ($n = 20$), compared to 5.33 for non-package cultures ($n = 75$)---not significantly different. This null result confirms that the \textit{structure} is independent of soil chemistry. Testing whether \textit{timing} correlates with soil chemistry requires interval data that Pulotu does not encode, representing a priority for future research.


%% ===================================================================
%% SECTION 5: TAPHONOMIC CALIBRATION
%% ===================================================================
\section{Taphonomic Calibration: The 1000-Day Hypothesis}

We now turn to the central question: what generated the specific numbers 3, 7, 40, 100, and 1000?

\subsection{Decomposition in Tropical Volcanic Soil}

The environmental conditions of lowland Java---mean annual temperature of 26--28\textdegree C, relative humidity above 70\%, annual rainfall exceeding 2,000~mm, and volcanic ash soils with pH 4.5--5.5---create a distinctive decomposition regime.

The fresh stage typically concludes within 1--3 days in tropical conditions \citep{rodriguez_bass_1985, galloway_1997}. The bloat stage reaches maximum intensity between days 3 and 7 \citep{galloway_1989, vass_2001}. Active decay progresses rapidly in warm, moist, acidic soil; at 28\textdegree C, 40 days corresponds to approximately 1,120 accumulated degree-days (ADD), consistently associated with advanced decay \citep{megyesi_2005}. By 100 days ($\sim$2,800 ADD), skeletonisation is typically well advanced \citep{haglund_sorg_1997}.

The critical variable for the 1000-day endpoint is bone preservation in acidic soil. Hydroxyapatite is chemically stable above pH~6.0 but undergoes progressive dissolution below this threshold \citep{nielsen_marsh_2000}. \citet{oghenemavwe_2022} demonstrated measurable bone degradation at pH~2.98 after six weeks. Archaeological case studies from Star Carr (pH 2--4) and Sutton Hoo (pH $\sim$4) document near-complete bone dissolution within years to decades.

At 1000 days ($\sim$2.7 years), significant bone degradation is taphonomically plausible in Javanese volcanic soil. While controlled experiments in andosol are lacking, the convergence of experimental short-term data, archaeological case studies, and hydroxyapatite dissolution chemistry supports this inference.

\subsection{The Interval Mapping}

The correspondence between slametan intervals and decomposition stages is not merely inferred from Western forensic science. Javanese tradition preserves explicit descriptions of the body's physical state at each milestone. The concept of \textit{kasampurnan} (perfection, completion) organises the sequence: each slametan marks the completion of decomposition for specific body components \citep{kalimullah_2016, soemodidjojo_1994}.

\begin{table}[H]
\centering
\caption{Slametan Intervals Mapped to Decomposition Stages}
\label{tab:mapping}
\small
\begin{tabularx}{\textwidth}{@{}ccXXc@{}}
\toprule
\textbf{Day} & \textbf{Name} & \textbf{Javanese Belief (Body State)} & \textbf{Forensic Stage (pH 4.5--5.5)} & \textbf{ADD} \\
\midrule
3 & Nelung dina & Nafsu dissipate; spirit in house & Fresh stage ends; autolysis complete & 84 \\
7 & Mitung dina & Skin and hair separate; spirit departs & Peak bloat; putrefactive gases maximal & 196 \\
40 & Matang puluh & Blood, flesh, marrow, innards, nails, hair, bones, muscle `perfected' & Advanced decay; soft tissue consumed & 1,120 \\
100 & Nyatus & Physical body (\textit{badan wadhag}) perfected & Skeletonisation advanced & 2,800 \\
730 & Mendhak pindo & `Body almost dissolved, only bones remain' & Bone surface degradation & 20,440 \\
1000 & Nyewu & `Body fully one with the earth' & Significant bone mineral dissolution & 28,000 \\
\bottomrule
\end{tabularx}
\end{table}

The slametan sequence tracks the progressive loss of body components in the exact order predicted by forensic taphonomy: vital processes $\rightarrow$ skin $\rightarrow$ soft tissue $\rightarrow$ skeleton $\rightarrow$ bone $\rightarrow$ earth. Figure~\ref{fig:mapping} presents this correspondence visually on a logarithmic timescale.

[Figure 1 about here]

\noindent\textit{Figure 1.} Slametan mortuary intervals mapped to forensic decomposition stage ranges. Each diamond marks a slametan ceremony; coloured bars show the range reported in forensic literature for each stage in tropical acidic conditions. All five intervals fall within their corresponding stage range.
\label{fig:mapping}

\subsection{Testing the correspondence: permutation and Monte Carlo}

A natural objection is that any set of five increasing numbers could be retrospectively mapped onto five decomposition stages---that the correspondence is an artefact of pattern-fitting. We test this objection with two complementary approaches, the first of which requires no researcher-defined parameters.

The strongest test is an exact permutation analysis. The five slametan numbers (3, 7, 40, 100, 1000) are fixed; we ask how many of the $5! = 120$ possible assignments of these numbers to the five decomposition stages produce a complete match. Since each stage has a defined position in the taphonomic sequence (fresh $\rightarrow$ bloat $\rightarrow$ advanced decay $\rightarrow$ skeletonisation $\rightarrow$ bone dissolution), we ask whether each slametan number falls within the forensic literature's reported range for its assigned stage. Only 1 of 120 permutations produces a complete match---the observed assignment---yielding $p = 0.008$. This test is robust to the choice of stage boundaries because the slametan numbers are sufficiently separated that few permutations place all five within plausible ranges regardless of how those ranges are defined.

As a supplementary analysis, we conducted a Monte Carlo simulation using stage boundaries derived from published forensic literature: fresh stage ends at 1--5 days \citep{rodriguez_bass_1985, galloway_1997}, bloat peaks at 3--14 days \citep{galloway_1989, vass_2001}, advanced decay at 20--80 days \citep{megyesi_2005}, skeletonisation advanced at 60--300 days \citep{haglund_sorg_1997}, and bone mineral dissolution significant at 300--2500 days \citep{nielsen_marsh_2000, oghenemavwe_2022}. We acknowledge that these boundaries are approximate and drawn from non-Javanese contexts; no published study exists for Javanese Andosol specifically. Drawing 500,000 random sets of five ordered positive integers from [1, 1500], only 11 (0.002\%) matched all five stages ($p < 0.001$). Under a conservative log-uniform distribution, $p = 0.009$ with narrow boundaries. The result is robust across a threefold sensitivity analysis varying stage widths from narrow to wide definitions, and across domain sizes from 500 to 3000 days.

We further tested six other documented mortuary traditions against the same stage boundaries (Table~\ref{tab:cross_tradition}). No other tradition achieved a complete five-stage match.

\begin{table}[H]
\centering
\caption{Cross-Tradition Decomposition Stage Matching}
\label{tab:cross_tradition}
\small
\begin{tabularx}{\textwidth}{@{}lXcc@{}}
\toprule
\textbf{Tradition} & \textbf{Mortuary intervals (days)} & \textbf{Stages matched} & \textbf{Soil context} \\
\midrule
\textbf{Javanese slametan} & \textbf{3, 7, 40, 100, 1000} & \textbf{5/5} & \textbf{Volcanic pH 4.5--5.5} \\
Toraja Rambu Solo & 1, 7, 30, 365, 730 & 4/5 & Volcanic pH 5.0 \\
Eastern Orthodox & 3, 9, 40, 365 & 4/4$^\dagger$ & Various \\
Buddhist bardo & 7, 14, 21, 28, 35, 42, 49, 100 & 3/5 & Various \\
Hindu shraddha & 3, 10, 13, 365 & 3/4$^\dagger$ & Various \\
Merina famadihana & 365, 730, 1095, 1825, 2555 & 1/5 & Laterite pH 4.7 \\
\bottomrule
\multicolumn{4}{l}{\footnotesize $^\dagger$ Fewer than 5 intervals; denominator is number of intervals.}
\end{tabularx}
\end{table}

Figure~\ref{fig:montecarlo} visualises the permutation and Monte Carlo results.

[Figure 2 about here]

\noindent\textit{Figure 2.} Statistical validation. (a) Exact permutation test: of 120 possible assignments of the five slametan numbers to five decomposition stages, only one produces a complete match ($p = 0.008$). (b) Monte Carlo simulation: of 200,000 random five-interval sets drawn uniformly from [1, 1500], only two achieve a 5/5 match ($p < 0.001$). Log scale on the $y$-axis emphasises the rarity of high-match outcomes.
\label{fig:montecarlo}

These results do not prove that the slametan intervals were derived from decomposition observation, but they demonstrate that the correspondence is unlikely to be coincidental. The fact that the Toraja---also in volcanic soil---achieve the next-highest match (4/5) is consistent with, though not proof of, the environmental calibration hypothesis. Individual intervals, particularly the widely attested $\sim$40-day mark, appear in other traditions, which may itself reflect a broader taphonomic pattern. The uniqueness claim here concerns not individual numbers but the complete five-interval system, and especially the 1000-day terminal ceremony, which is currently unattested outside Java in our comparative sample.

\subsection{The observation mechanism}

A second objection asks how Javanese communities could have accumulated knowledge of underground decomposition without practising exhumation. We suggest three complementary pathways, ordered by evidential strength.

The primary pathway is olfactory. Putrefactive gases produced during the bloat stage (days 3--14) are well documented in the forensic literature as detectable through soil, particularly in shallow burials \citep{vass_2001, carter_tibbett_2008}. The Primbon specifies burial at approximately 75~cm with a protective cover (\textit{glogor}) but direct earth contact (Entry No.~333). At this depth in porous volcanic soil, the intense odour of the bloat phase would be perceptible at the surface. The olfactory transition---from strong smell (days 3--7) to fading smell (days 7--40) to no smell (after day 40)---provides a coarse but unmistakable signal of the body's transformation, requiring no disturbance of the grave. This pathway alone could account for the first three slametan intervals.

The secondary pathway is incidental exposure. In densely used burial grounds, the digging of new graves periodically exposes older remains at various stages of decomposition. Flooding, erosion, and animal activity---all common in tropical lowland Java---would occasionally reveal buried remains. Over generations, such incidental exposure would accumulate into a practical knowledge base of decomposition progression, even without deliberate exhumation.

The third pathway is institutionalised grave visitation. The Primbon prescribes visitation during \textit{bulan Ruwah} (Entry No.~334), during which offerings are placed at the gravestone. Regular visitation creates opportunities to observe the general condition of the grave site---whether the ground has settled, whether vegetation has changed, whether the grave mound has collapsed---though we acknowledge that such observations are coarse indicators, not precise measurements.

We do not claim that any single pathway is sufficient. Rather, the combination of olfactory signals, incidental exposure, and periodic visitation, sustained over many generations in a community practising consistent shallow burial in the same soil environment, provides a plausible---though not proven---mechanism by which empirical knowledge of decomposition timing could have accumulated.

\subsection{The pH--Timing Prediction}

If the slametan intervals reflect empirical observation of decomposition in volcanic soil, a testable prediction follows: communities in alkaline soils (pH~7--8), where bone mineral is stable for centuries, should either lack a fixed terminal ceremony analogous to \textit{nyewu} or set this ceremony at a substantially later date. The Berawan of Borneo (soil pH~4.8), who store corpses in above-ground containers to observe decomposition directly before secondary burial \citep{metcalf_1982}, represent a particularly instructive case: their practice of monitoring decomposition as a ritual guide is precisely the behaviour the hypothesis predicts.

\subsection{Caveats}

We do not claim exact correspondence. Decomposition rates are highly variable, and no published study has been conducted specifically in Javanese andosol. Our claim is that the \textit{ordering} and \textit{approximate magnitude} are consistent with generations of empirical observation in acidic volcanic soil. The progression from odour through visible decay to bare bone to dissolution into earth is observable to any community practising shallow burial and periodic grave visitation.


%% ===================================================================
%% SECTION 6: H-TOM SYNTHESIS
%% ===================================================================
\section{The H-TOM Synthesis: Two Timescales, One Driver}

\subsection{The Ritual Timescale: 1000 Days}

The slametan cycle operates on what we term the ritual timescale. Over approximately 1000 days, the body of the deceased undergoes progressive dissolution in acidic volcanic soil. The slametan marks this progression at empirically significant thresholds. At the endpoint---\textit{nyewu}---the body has `completely decayed to dust' \citep[ch.~6]{geertz_1960}, and the soul `will no longer return to the family' \citep[440]{hendrajaya_2020}.

The Primbon Betaljemur Adammakna codifies this calendar in computational form. Entry No.~332 presents a calculation table mapping the day and \textit{pasaran} of death to the calendrical dates for each subsequent ceremony through \textit{nyewu} \citep[245--246]{soemodidjojo_1994}. The entry defines \textit{nyewu} explicitly as `selamatan hari keseribu terhitung dari meninggal dunianya si mati'---the ceremony on the thousandth day counted from the death of the deceased. That this sequence is treated as a fixed computation implies that the intervals were understood as natural and necessary.

Immediately preceding this table in the Primbon is Entry No.~331, a prayer for cases in which `the corpse smells bad' (\textit{jika jenasah berbau busuk}). The juxtaposition of a decomposition-management entry with the selamatan calculation table reflects an indigenous framework in which ritual timing and bodily decay are conceptually contiguous.

\subsection{The Geological Timescale: 1000 Years}

A second, slower taphonomic process operates on the same landscape. Archaeological evidence from multiple sites documents composite sedimentation rates of 3.5--7~mm/yr in Java's volcanic plains. Candi Sambisari (6.5~m burial over $\sim$1000 years), Candi Kedulan (7~m burial), the Liangan settlement near Mt.\ Sindoro (5--10~m burial), and the Dwarapala of Singhasari (185~cm in 535 years) demonstrate the scale of this process \citep{newhall_2000}. At these rates, a surface site is buried beneath several metres---beyond casual discovery---within $\sim$1000 years. These are monumental stone structures; domestic wooden architecture fares far worse.

\subsection{The Convergence}

The critical insight is that these two timescales are driven by a single volcanic process. The tephra that weathers into acidic Andosol is the same material that accumulates as the sedimentary matrix burying sites. Volcanism simultaneously creates the soil chemistry that dissolves bodies rapidly and produces the depositional environment that removes archaeological sites from visibility.

\begin{table}[H]
\centering
\caption{Two Timescales of Volcanic Taphonomic Bias}
\label{tab:timescales}
\begin{tabularx}{\textwidth}{@{}lXX@{}}
\toprule
\textbf{Parameter} & \textbf{Ritual Timescale} & \textbf{Geological Timescale} \\
\midrule
Duration & $\sim$1,000 days ($\sim$2.7 years) & $\sim$1,000 years \\
Process & Body decomposition in acidic soil & Site burial by volcanic sedimentation \\
Driver & Tephra-derived Andosol (pH 4.5--5.5) & Tephra deposition ($\sim$3.5 mm/yr) \\
Endpoint & Body dissolved (\textit{nyewu}) & Site buried $>$3.5 m \\
Cultural response & Slametan ritual calendar & None (invisible to single generation) \\
\bottomrule
\end{tabularx}
\end{table}

This generates what we term the Hypothetical Taphonomic Overlay Model (H-TOM): ritual and geological taphonomy as complementary expressions of the same environmental forcing (Figure~\ref{fig:htom}). The living culture records, in ritual, the very process that erases its own material past.

[Figure 3 about here]

\noindent\textit{Figure 3.} The Hypothetical Taphonomic Overlay Model (H-TOM). Volcanism drives two taphonomic processes at different timescales: chemical weathering of tephra-derived Andosol dissolves buried remains within $\sim$1,000 days, while tephra deposition buries archaeological sites within $\sim$1,000 years. The slametan ritual calendar tracks the first process; both processes are invisible to any single generation.
\label{fig:htom}


%% ===================================================================
%% SECTION 7: DISCUSSION
%% ===================================================================
\section{Discussion}

\subsection{Limitations}

Several limitations constrain our claims. First, the decomposition timelines are drawn from general forensic literature, not site-specific Javanese cemetery data. No published decomposition study has been conducted in Javanese volcanic Andosol---the most significant gap in the evidentiary chain.

Second, the slametan intervals as documented today may not be identical to their historical antecedents. However, three observations mitigate this concern: (a) multiple scholars independently classify the intervals as pre-dating both Islamic and Hindu-Buddhist influence; (b) the intervals are consistent across geographically separated communities; and (c) the Primbon treats them as a fixed computational system.

Third, Pulotu codes beliefs at the culture level; within-culture variation is not captured. Fourth, ISRIC soil data represent modern conditions, though the primary determinant of Andosol acidity---volcanic parent material mineralogy---is geologically stable on millennial timescales.

\subsection{Falsification Criteria}

\textbf{Prediction 1 (Strong test):} Controlled decomposition experiments in Javanese Andosol should produce timing broadly consistent with the slametan intervals. If such conditions produce a radically different trajectory, the hypothesis is falsified.

\textbf{Prediction 2 (Cross-cultural test):} Communities in alkaline soils (pH~7--8) should either lack a terminal ceremony analogous to \textit{nyewu} or set it at a substantially later date. If alkaline-soil cultures exhibit the same 1000-day endpoint, the environmental calibration hypothesis is undermined.

\textbf{Prediction 3 (Primbon test):} If comprehensive textual analysis reveals an explicit theological rationale for the intervals, independent of bodily state, the hypothesis would need modification.

\textbf{Prediction 4 (Grave visitation test):} If ethnographic fieldwork demonstrates that Javanese communities did not practise periodic grave visitation during the slametan cycle, the observational mechanism would be weakened. The Primbon's inclusion of a grave visitation protocol (Entry No.~334) suggests such visitation was institutionalised.

\subsection{The `Just-So Story' Objection}

We address this objection in four ways. First, the hypothesis was generated from the data: the correspondence between slametan intervals and decomposition stages emerged from independent comparison of forensic and ethnographic sources. Second, it generates specific, falsifiable predictions (Section 7.2). Third, we do not claim exact correspondence---only that the ordering and approximate magnitude are consistent with empirical observation. Fourth, the permutation test reported in Section 5---which requires no researcher-defined stage boundaries---yields $p = 0.008$ for the observed interval-to-stage assignment. Supplementary Monte Carlo simulations confirm this result across a range of parameter settings. The correspondence is quantifiably unlikely to be coincidental.

\subsection{The Diffusion Alternative}

An alternative explanation for the 40-day interval is cultural diffusion. The Sufi practice of \textit{chilla} (40-day spiritual retreat), well established in the Islamic networks that reached Java from the 13th century onward, provides a plausible transmission channel for the number 40 \citep{woodward_1989}. \citet{woodward_1989} has argued that the slametan is fundamentally Islamic/Sufi in character, with its forms originating from Islamic texts rather than pre-Islamic tradition.

We acknowledge this possibility for the individual number 40, but note three constraints on the diffusion argument. First, the Hindu shraddha system---the pre-Islamic death ritual framework in Java---contains no 40-day interval (its cycle runs at 10--13 days). If 40 entered Javanese mortuary practice through Islam, it entered a system that already possessed the other intervals (3, 7, 100, 1000) independently. Second, the Sufi chilla is a practice for the living, not a mortuary observance; no Islamic or Sufi source prescribes a 1000-day death commemoration. Third, and most importantly, the complete sequence 3--7--40--100--1000 is unattested in any non-Javanese tradition. Even if 40 was reinforced through contact with Sufi Islam, the system as a whole---and particularly its terminal \textit{nyewu}---remains a Javanese formation.

Indeed, the cross-cultural recurrence of a $\sim$40-day mortuary interval may itself support the taphonomic hypothesis rather than undermine it. Ancient Egyptian mummification required 40 days of natron treatment---a duration empirically calibrated to the time needed to desiccate a human body. If multiple unconnected cultures independently converged on approximately 40 days as a mortuary milestone, the most parsimonious explanation is that they were all observing the same biological process: the approximate completion of soft tissue decomposition.

\subsection{Broader Significance}

The hypothesis suggests that mortuary ritual timing may constitute a previously unrecognised source of environmental information---a form of `taphonomic ethnography' in which living cultural practices encode empirical knowledge about local decomposition conditions. It also contributes to debates about the relationship between `belief' and `observation' in non-Western knowledge systems. The slametan tradition's language---its inventory of tissue types, its claims about bodily states at specific intervals, its computational calendar---is observational in character, even as its practitioners classify it as religious ritual.


%% ===================================================================
%% SECTION 8: CONCLUSION
%% ===================================================================
\section{Conclusion}

This paper has advanced three principal findings.

First, the Javanese slametan death sequence---commemorations at 3, 7, 40, 100, and 1000 days after death---cannot be derived from Hindu, Buddhist, Islamic, or Christian scripture. The intervals are absent from shraddha liturgy, inconsistent with bardo cosmology, explicitly classified as \textit{bid'ah} by Islamic jurisprudence, and have no precedent in Christian tradition. That four successive religious waves---Hinduism (from $\sim$4th c.\ CE), Buddhism, Islam (from $\sim$13th c.\ CE), and Christianity (from $\sim$16th c.\ CE)---each modified the verbal content of the slametan while preserving its temporal structure confirms the intervals belong to the indigenous Javanese substrate. The epigraphic record of 268 Nusantaran inscriptions further confirms the existence of an indigenous cosmological vocabulary that persisted throughout the period of Indianisation.

Second, the slametan intervals correspond to the stages of human decomposition predicted by forensic taphonomy for a body buried in acidic volcanic soil. A permutation test demonstrates that the probability of this correspondence arising by chance is $p = 0.008$, and no other mortuary tradition in our comparative sample achieves a complete five-stage match. The Javanese tradition's own descriptions---the tissue inventory at 40 days, the `only bones remain' at two years, the `fully one with the earth' at 1000 days---match the taphonomic predictions with a specificity suggesting observational knowledge. The Primbon Betaljemur Adammakna reinforces this inference by codifying the schedule as a fixed computation and placing a decomposition-management entry immediately before the calculation table. We propose that olfactory signals from shallow burial, combined with incidental exposure and periodic grave visitation, provided a plausible mechanism for the accumulation of this knowledge without deliberate exhumation.

Third, the structural framework---death as gradual process, post-mortem ritual efficacy, the soul's journey paralleling the body's decay---is pan-Austronesian, shared by at least 30 of 137 cultures in the Pulotu database. The structure is inherited; the specific intervals are locally variable, consistent with environmental calibration.

We have proposed the `Volcanic Ritual Clock' hypothesis: that the slametan intervals were empirically calibrated, through generations of observation, to the pace of decomposition in volcanic soil. We do not reduce the slametan to a decomposition calendar. Its social, emotional, and cosmological dimensions are real and irreducible. We argue that the physical process of decomposition provided the empirical scaffold upon which those richer meanings were built---and that this scaffold is still legible in the tradition's own language about the body.

The Volcanic Ritual Clock operates at two timescales linked by a single volcanic driver. At the ritual timescale, approximately 1000 days, acidic volcanic soil reduces a buried body to dust. At the geological timescale, approximately 1000 years, volcanic sedimentation buries the archaeological sites of the communities that performed these rituals. The environment that calibrates the ritual clock also destroys the material evidence of the people who kept it. The living culture records, in ritual, the very process that erases its own material past.


%% ===================================================================
%% AUTHOR NOTE — REMOVED FOR DOUBLE-BLIND REVIEW
%% Restore after acceptance.
%% ===================================================================

% ANONYMIZED: Author Positionality and AI Disclosure removed for review.
% See commented block below for post-acceptance restoration.

\subsection*{Methodological Note}

The contribution claimed here is methodological, not ethnographic. This paper brings computational tools---corpus analysis, Monte Carlo simulation, cross-database querying---to a question that has been posed qualitatively by anthropologists (Hertz, Bloch, Metcalf) but never tested quantitatively. The interpretive claims about Javanese culture draw entirely on published ethnographic and textual sources; no original fieldwork was conducted. This research was conducted with the assistance of a large language model (LLM) for systematic literature review, statistical computation, and corpus analysis. All hypotheses, interpretations, and scientific claims are the author's own.

% === POST-ACCEPTANCE: RESTORE THE FOLLOWING ===
%\subsection*{Author Positionality}
%
%The author is a data scientist at Universitas Bhinneka Nusantara, Malang (East Java),
%with training in computer science and machine learning, not in archaeology or anthropology.
%This paper grew from a personal observation: that Javanese communities in the author's
%home region continue to practise the slametan death cycle with precise numerical intervals,
%while the archaeological record of those same communities is systematically absent from
%the landscape. The question of \textit{why these numbers} arose not from fieldwork but
%from computation---from noticing that the intervals matched a forensic decomposition
%trajectory when parameterised for local soil conditions.
%
%The contribution claimed here is methodological, not ethnographic. The author brings
%computational tools---corpus analysis, Monte Carlo simulation, cross-database querying---to
%a question that has been posed qualitatively by anthropologists (Hertz, Bloch, Metcalf)
%but never tested quantitatively. The interpretive claims about Javanese culture draw
%entirely on published ethnographic and textual sources; no original fieldwork was conducted.
%The author hopes this computational framework may be useful to Javanese archaeologists
%and anthropologists who share the concern---expressed by Stutterheim (1948) and others---that
%Nusantaran civilisational depth is obscured by taphonomic bias rather than genuinely absent.
%
%\subsection*{AI Disclosure}
%
%This research was conducted with the assistance of Claude (Anthropic), a large language
%model, used for: systematic literature review across Indonesian-, English-, and
%Dutch-language sources; corpus analysis of the DHARMA epigraphic database; statistical
%computation (Monte Carlo simulation, permutation testing); cross-database querying of
%the Pulotu ethnographic dataset; and iterative drafting. All hypotheses, interpretations,
%and scientific claims are the author's own. The computational workflow is documented in
%the project repository.
% === END POST-ACCEPTANCE BLOCK ===


%% ===================================================================
%% ETHICS STATEMENT
%% ===================================================================
\subsection*{Ethics Statement}

This study involved no human participants, fieldwork, or collection of personal data. All data sources are publicly available: the DHARMA epigraphic corpus (CC-BY 4.0), the Pulotu ethnographic database (CC-BY 4.0), ISRIC SoilGrids (open access), and published forensic taphonomy literature. Ethnographic descriptions of mortuary practices are drawn from published scholarly sources.

%% ===================================================================
%% REFERENCES
%% ===================================================================
\newpage
\begin{thebibliography}{99}

\bibitem[Adams(2006)]{adams_2006}
Adams, Kathleen M.\ (2006). \textit{Art as politics: Re-crafting identities, tourism, and power in Tana Toraja, Indonesia}. Honolulu: University of Hawai'i Press.

\bibitem[Adelaar(1995)]{adelaar_1995}
Adelaar, K.\ Alexander (1995). `Asian roots of the Malagasy: A linguistic perspective', \textit{Bijdragen tot de Taal-, Land- en Volkenkunde} 151-3:325--56.

\bibitem[Aizid(2015)]{aizid_2015}
Aizid, Rizem (2015). \textit{Islam abangan dan kehidupannya: Seluk-beluk kehidupan Islam abangan}. Yogyakarta: Dipta.

\bibitem[Aritonang and Steenbrink(2008)]{aritonang_steenbrink_2008}
Aritonang, Jan Sihar and Karel Steenbrink (2008). \textit{A history of Christianity in Indonesia}. Leiden: Brill.

\bibitem[Beatty(1999)]{beatty_1999}
Beatty, Andrew (1999). \textit{Varieties of Javanese religion: An anthropological account}. Cambridge: Cambridge University Press.

\bibitem[Bloch(1971)]{bloch_1971}
Bloch, Maurice (1971). \textit{Placing the dead: Tombs, ancestral villages, and kinship organization in Madagascar}. London: Seminar Press.

\bibitem[Bloch and Parry(1982)]{bloch_parry_1982}
Bloch, Maurice and Jonathan Parry (1982). \textit{Death and the regeneration of life}. Cambridge: Cambridge University Press.

\bibitem[Blust and Trussel(2010)]{blust_trussel_2010}
Blust, Robert and Stephen Trussel (2010). \textit{Austronesian Comparative Dictionary}.\footnote{\url{www.trussel2.com/ACD} (accessed 9-3-2026). Online, ongoing.}

\bibitem[Bruce-Mitford(1975)]{bruce_mitford_1975}
Bruce-Mitford, Rupert (1975). \textit{The Sutton Hoo ship-burial}. London: British Museum.

\bibitem[Carter and Tibbett(2008)]{carter_tibbett_2008}
Carter, David O.\ and Mark Tibbett (2008). `Cadaver decomposition and soil: Processes', in: \textit{Soil analysis in forensic taphonomy}, pp.\ 29--52. Boca Raton: CRC Press.

\bibitem[Cuevas(2003)]{cuevas_2003}
Cuevas, Bryan J.\ (2003). \textit{The hidden history of the Tibetan Book of the Dead}. Oxford: Oxford University Press.

\bibitem[Galloway(1997)]{galloway_1997}
Galloway, Alison (1997). `The process of decomposition: A model from the Arizona-Sonoran desert', in: William D.\ Haglund and Marcella H.\ Sorg (eds), \textit{Forensic taphonomy: The postmortem fate of human remains}, pp.\ 139--50. Boca Raton: CRC Press.

\bibitem[Galloway et~al.(1989)]{galloway_1989}
Galloway, Alison et~al.\ (1989). `Decay rates of human remains in an arid environment', \textit{Journal of Forensic Sciences} 34-3:607--16.

\bibitem[Gammage(2011)]{gammage_2011}
Gammage, Bill (2011). \textit{The biggest estate on earth: How Aborigines made Australia}. Sydney: Allen \& Unwin.

\bibitem[Geertz(1960)]{geertz_1960}
Geertz, Clifford (1960). \textit{The religion of Java}. Chicago: University of Chicago Press.

\bibitem[Graeber(1995)]{graeber_1995}
Graeber, David (1995). `Dancing with corpses reconsidered: An interpretation of \textit{famadihana}', \textit{American Ethnologist} 22-2:258--78.

\bibitem[Griffiths et~al.(2020)]{griffiths_2020}
Griffiths, Arlo et~al.\ (2020). \textit{DHARMA: The Domestication of ``Hindu'' Asceticism and the Religious Making of South and Southeast Asia}. [ERC project, epigraphic corpus. CC-BY 4.0.]

\bibitem[Haglund and Sorg(1997)]{haglund_sorg_1997}
Haglund, William D.\ and Marcella H.\ Sorg (1997). \textit{Forensic taphonomy: The postmortem fate of human remains}. Boca Raton: CRC Press.

\bibitem[Haglund and Sorg(2002)]{haglund_sorg_2002}
Haglund, William D.\ and Marcella H.\ Sorg (2002). \textit{Advances in forensic taphonomy: Method, theory, and archaeological perspectives}. Boca Raton: CRC Press.

\bibitem[Hendrajaya and Almu'tasim(2020)]{hendrajaya_2020}
Hendrajaya, Agus and Jumaidah Almu'tasim (2020). `Tradisi selamatan kematian nyatus nyewu: Implikasi nilai pluralisme Islam Jawa', \textit{Jurnal Lektur Keagamaan} 17-2:431--60.

\bibitem[Hertz(1960)]{hertz_1907}
Hertz, Robert (1960). `A contribution to the study of the collective representation of death', in: \textit{Death and the right hand}, pp.\ 27--86. London: Cohen \& West. [Original published 1907.]

\bibitem[Junker(1999)]{junker_1999}
Junker, Laura Lee (1999). \textit{Raiding, trading, and feasting: The political economy of Philippine chiefdoms}. Honolulu: University of Hawai'i Press.

\bibitem[Kalimullah(2016)]{kalimullah_2016}
Kalimullah, M.\ (2016). \textit{Primbon dalam budaya Jawa}. [PhD thesis, UIN Sunan Ampel Surabaya.]

\bibitem[Kemp(2016)]{kemp_2016}
Kemp, W.L.\ (2016). `High soil acidity associated with near complete mineral dissolution of recently buried human remains'. [ResearchGate.]

\bibitem[Knipe(1977)]{knipe_1977}
Knipe, David M.\ (1977). `Sapindikarana: The Hindu rite of entry into heaven', in: Frank E.\ Reynolds and Earle H.\ Waugh (eds), \textit{Religious encounters with death}, pp.\ 111--24. University Park: Penn State University Press.

\bibitem[Koentjaraningrat(1994)]{koentjaraningrat_1994}
Koentjaraningrat (1994). \textit{Kebudayaan Jawa}. Jakarta: Balai Pustaka.

\bibitem[Megyesi, Nawrocki and Haskell(2005)]{megyesi_2005}
Megyesi, Mary S., Stephen P.\ Nawrocki and Neal H.\ Haskell (2005). `Using accumulated degree-days to estimate the postmortem interval from decomposed human remains', \textit{Journal of Forensic Sciences} 50-3:618--26.

\bibitem[Metcalf(1982)]{metcalf_1982}
Metcalf, Peter (1982). \textit{A Borneo journey into death: Berawan eschatology from its rituals}. Philadelphia: University of Pennsylvania Press.

\bibitem[Metcalf and Huntington(1991)]{metcalf_huntington_1991}
Metcalf, Peter and Richard Huntington (1991). \textit{Celebrations of death: The anthropology of mortuary ritual}. Second edition. Cambridge: Cambridge University Press.

\bibitem[Milner et~al.(2016)]{milner_2016}
Milner, Nicky et~al.\ (2016). `Lessons from Star Carr on the study of Mesolithic environments', \textit{Proceedings of the National Academy of Sciences} 113-46:13182--7.

\bibitem[Newhall et~al.(2000)]{newhall_2000}
Newhall, Christopher G.\ et~al.\ (2000). `10,000 years of explosive eruptions of Merapi volcano, Central Java: Archaeological and modern implications', \textit{Journal of Volcanology and Geothermal Research} 100-1/4:9--50.

\bibitem[Nielsen-Marsh and Hedges(2000)]{nielsen_marsh_2000}
Nielsen-Marsh, Christina M.\ and Robert E.M.\ Hedges (2000). `Patterns of diagenesis in bone I: The effects of site environments', \textit{Journal of Archaeological Science} 27:1139--50.

\bibitem[Nooy-Palm(1979)]{nooy_palm_1979}
Nooy-Palm, Hetty (1979). \textit{The Sa'dan-Toraja: A study of their social life and religion}. Vol.\ 1. The Hague: KITLV Press. [Verhandelingen KITLV 87.]

\bibitem[Oghenemavwe et~al.(2022)]{oghenemavwe_2022}
Oghenemavwe, E.L.\ et~al.\ (2022). `Histomorphometric analysis of buried bone in acidic soil', \textit{Annals of Biomedical and Health Sciences}.

\bibitem[Parry(1994)]{parry_1994}
Parry, Jonathan (1994). \textit{Death in Banaras}. Cambridge: Cambridge University Press.

\bibitem[Poggio et~al.(2021)]{poggio_2021}
Poggio, Laura et~al.\ (2021). `SoilGrids 2.0: Producing soil information for the globe with quantified spatial uncertainty', \textit{SOIL} 7:217--40.

\bibitem[Rodriguez and Bass(1985)]{rodriguez_bass_1985}
Rodriguez, William C.\ and William M.\ Bass (1985). `Decomposition of buried bodies and methods that may aid in their location', \textit{Journal of Forensic Sciences} 30-3:836--52.

\bibitem[Sholihah et~al.(2023)]{sholihah_2023}
Sholihah, A.W.\ et~al.\ (2023). `Tinjauan filosofis tradisi selamatan orang meninggal di Jawa dalam perspektif Islam', in: \textit{Proceedings of the 2nd ICCL}, pp.\ 360--73. Surakarta: UIN Raden Mas Said Surakarta.

\bibitem[Soemodidjojo(1994)]{soemodidjojo_1994}
Soemodidjojo, M.\ (1994). \textit{Kitab Primbon Betaljemur Adammakna}. Yogyakarta: Soemodidjojo Mahadewa / CV.\ Buana Raya. [Indonesian translation by Ir.\ Wibatsu Harianto. Original in Javanese.]

\bibitem[Suwito, Sriyanto and Hidayat(2015)]{suwito_2015}
Suwito, Sriyanto and Hidayat (2015). `Tradisi dan ritual kematian wong Islam Jawa', \textit{IBDA': Jurnal Kebudayaan Islam} 13-2:197--216.

\bibitem[Teiser(1988)]{teiser_1988}
Teiser, Stephen F.\ (1988). \textit{The ghost festival in medieval China}. Princeton: Princeton University Press.

\bibitem[Vass(2001)]{vass_2001}
Vass, Arpad A.\ (2001). `Beyond the grave: Understanding human decomposition', \textit{Microbiology Today} 28:190--3.

\bibitem[Volkman(1985)]{volkman_1985}
Volkman, Toby Alice (1985). \textit{Feasts of honor: Ritual and change in the Toraja highlands}. Urbana: University of Illinois Press.

\bibitem[Watts et~al.(2015)]{watts_2015}
Watts, Joseph et~al.\ (2015). `Pulotu: Database of Austronesian supernatural beliefs and practices', \textit{PLOS ONE} 10-9:e0136783.

\bibitem[Woodward(1989)]{woodward_1989}
Woodward, Mark R.\ (1989). \textit{Islam in Java: Normative piety and mysticism in the Sultanate of Yogyakarta}. Tucson: University of Arizona Press.

\end{thebibliography}

\end{document}
